% ================================================================
%  БҮЛЭГ 2: СИСТЕМИЙН ШААРДЛАГА БА ШИНЖИЛГЭЭ
% ================================================================
\chapter{СИСТЕМИЙН ШААРДЛАГА БА ШИНЖИЛГЭЭ}
\label{ch2}

\section{Системийн үйл ажиллагааны тухай дэлгэрэнгүй}

Санал болгож буй систем нь дараах бүрэн урсгалыг хэрэгжүүлнэ (Зураг~\ref{fig:flow}):

\begin{enumerate}[leftmargin=*]
  \item Үйлчлүүлэгч ресторанд орж ширээнд суухад ширээн дэлгэц дээр QR код байрлана;
  \item Хэрэглэгч смартфоны камераар QR кодыг уншуулна;
  \item Мобайл аппликейшн автоматаар нээгдэж тухайн рестораны дижитал цэс 3 секундын
        дотор харагдана;
  \item Хэрэглэгч хоол сонгож сагсанд нэмэн ``Захиалга өгөх'' товчийг дарна;
  \item Захиалга нь WebSocket протоколоор ажилтны дэлгэцэнд бодит цаг хугацаанд дамжина;
  \item Цайчин захиалгыг баталгаажуулахад зочны дэлгэцэнд мэдэгдэл ирнэ;
  \item Тогооч хоол бэлтгэж ``Бэлэн боллоо'' тэмдэглэхэд зочны дэлгэцэнд мэдэгдэл ирнэ;
  \item Зочин тооцоо хийлгэхэд кассир тооцоог баталгаажуулна.
\end{enumerate}

\begin{figure}[htbp]
\centering
\begin{tikzpicture}[
  step/.style={rectangle, draw, rounded corners=4pt, minimum width=3.5cm,
    minimum height=0.9cm, align=center, font=\small, fill=blue!10},
  decision/.style={diamond, draw, aspect=2, minimum width=3cm,
    minimum height=0.7cm, align=center, font=\scriptsize, fill=yellow!20},
  arr/.style={-{Stealth}, thick},
  node distance=1.2cm
]
\node[step] (qr) {QR код уншуулах};
\node[step, below=of qr] (menu) {Цэс харуулах};
\node[step, below=of menu] (order) {Захиалга үүсгэх};
\node[step, below=of order] (ws) {WebSocket дамжуулалт};
\node[step, below=of ws] (confirm) {Ажилтан баталгаажуулах};
\node[step, below=of confirm] (prepare) {Гал тогоо бэлтгэх};
\node[step, below=of prepare] (serve) {Хоол хүргэх};
\node[step, below=of serve] (pay) {Тооцоо хийх};

\draw[arr] (qr) -- (menu);
\draw[arr] (menu) -- (order);
\draw[arr] (order) -- (ws);
\draw[arr] (ws) -- (confirm);
\draw[arr] (confirm) -- (prepare);
\draw[arr] (prepare) -- (serve);
\draw[arr] (serve) -- (pay);

\node[right=1.5cm of qr, font=\scriptsize, text=gray] {GET /api/tables/:token};
\node[right=1.5cm of order, font=\scriptsize, text=gray] {POST /api/orders};
\node[right=1.5cm of ws, font=\scriptsize, text=gray] {socket.emit('new\_order')};
\node[right=1.5cm of pay, font=\scriptsize, text=gray] {PATCH /api/orders/:id};
\end{tikzpicture}
\caption{Системийн ерөнхий үйл ажиллагааны урсгал}
\label{fig:flow}
\end{figure}

\section{Системийг ашиглах хэрэглэгчид}

Систем нь 3 үндсэн хэрэглэгчийн үүрэгтэй:

\begin{table}[htbp]
\caption{Системийн хэрэглэгчдийн үүргийн тодорхойлолт}
\label{tab:roles}
\centering
\small
\begin{tabularx}{\textwidth}{|L|L|L|}
\hline
\rowcolor{TableHdr}
\textcolor{white}{\textbf{Хэрэглэгч}} &
\textcolor{white}{\textbf{Хандах нөөц}} &
\textcolor{white}{\textbf{Боломжит үйлдэл}} \\ \hline
Зочин (Guest) & Цэс, захиалга &
  QR уншуулах, хоол сонгох, захиалга өгөх, статус харах \\ \hline
\rowcolor{LightBlue}
Кассир (Cashier) & Захиалга, тооцоо, нэгтгэл &
  Захиалга харах, тооцоо баталгаажуулах, тайлан харах \\ \hline
Менежер (Manager) & Бүх нөөц &
  Цэс/ширээ/захиалга/тайлан бүрэн удирдлага \\ \hline
\end{tabularx}
\end{table}

\section{Функционал шаардлага}

\textbf{ФШ-01: QR Кодоор Цэс Нээх}\\
Хэрэглэгч ширээн дэлгэц дээрх QR кодыг смартфоны камераар уншуулна. 3 секундын
дотор тухайн ширээний мэдээлэл болон ресторан цэс нэн даруй нээгдэнэ. Цэс нь
ангиллаар бүлэглэгдсэн, хоол бүрийн зураг, нэр, тайлбар, үнэ агуулсан байна.

\textbf{ФШ-02: Захиалга Үүсгэх}\\
Хэрэглэгч цэснээс хоол сонгон сагсанд нэмэх, хасах боломжтой байна. Хоол бүрд
тусгай хүсэлт бичих боломжтой байна. ``Захиалга өгөх'' товч дарахад нийт дүн болон
захиалгын мэдээллийн баталгаажуулалтын дэлгэц харагдана.

\textbf{ФШ-03: Бодит Цагийн Мэдэгдэл (Зочин)}\\
Захиалга баталгаажсан, бэлтгэж байна, бэлэн боллоо, хүргэгдсэн зэрэг статус
өөрчлөлт бүрийг зочны дэлгэцэнд шуурхай харуулна.

\textbf{ФШ-04: Бодит Цагийн Мэдэгдэл (Ажилтан)}\\
Шинэ захиалга ирэхэд ажилтны дэлгэцэнд дохиолол болон мэдэгдэл харагдана.
Захиалга нь ширээний нэр, хоолны жагсаалт, нийт дүн агуулсан байна.

\textbf{ФШ-05: Захиалгын Статус Удирдлага}\\
Захиалгын статусын шилжилт нь зөвхөн тодорхой дарааллаар явагдана:
\[
\texttt{pending} \to \texttt{confirmed} \to \texttt{preparing} \to
\texttt{ready} \to \texttt{served} \to \texttt{paid}
\]
Статус өөрчлөлт бүр WebSocket-р зочны дэлгэцэнд шуурхай мэдэгдэнэ.

\textbf{ФШ-06: Цэс Удирдлага}\\
Менежер цэсийн ангилал болон хоолыг нэмэх, засварлах, устгах боломжтой байна.
Хоолны байдал (байна/байхгүй) бодит цаг хугацаанд шинэчлэх боломжтой байна.

\textbf{ФШ-07: Ширээ Удирдлага}\\
Менежер ширээ нэмэх, засах, QR кодыг хэвлэх боломжтой байна. QR token
автоматаар UUID форматаар үүснэ.

\textbf{ФШ-08: Тайлан ба Нэгтгэл}\\
Менежер/кассир өдрийн орлого, ширээ тус бүрийн захиалгын тоо, хамгийн эрэлттэй
хоолны жагсаалтыг харах боломжтой байна. Тогтоосон хугацааны тайлан гаргана.

\section{Функционал бус шаардлага}

\subsection{Програм хангамжийн чанарын шаардлага}

\begin{table}[htbp]
\caption{Функционал бус шаардлагын хүснэгт}
\label{tab:nfr}
\centering
\small
\begin{tabularx}{\textwidth}{|L|L|Y|}
\hline
\rowcolor{TableHdr}
\textcolor{white}{\textbf{Шаардлага}} &
\textcolor{white}{\textbf{Хэмжүүр}} &
\textcolor{white}{\textbf{Заалт}} \\ \hline
Хурд & API хариу & 200 мс-ээс бага \\ \hline
\rowcolor{LightBlue}
Хурд & WS мессеж & 100 мс-ээс бага \\ \hline
Хүртээмж & Uptime & 99.5\% буюу түүнээс дээш \\ \hline
\rowcolor{LightBlue}
Ачаалал & Нэгэн цагт & 200+ холболт \\ \hline
Аюулгүй байдал & Нотлох & JWT + HTTPS/TLS \\ \hline
\rowcolor{LightBlue}
Хэрэглэхэд хялбар & SUS оноо & 80+ (Сайн) \\ \hline
Дэлгэцэнд зохицох & Responsive & 320px -- 1920px \\ \hline
\rowcolor{LightBlue}
Сонсогдох & WCAG & 2.1 AA түвшин \\ \hline
\end{tabularx}
\end{table}

\section{Ашигласан технологи, техникийн шаардлага}

\subsection{Технологийн шаардлага}

\textbf{Програм ажиллах үйлдлийн системийн сонголт}

Сервер талд Ubuntu 22.04 LTS үйлдлийн системийг ашиглана. Ubuntu нь бүх нийтийн
серверийн стандарт үйлдлийн систем бөгөөд Node.js-тэй нэвтрэлт сайтай.

\textbf{Програм хөгжүүлэх хэлний сонголт}

TypeScript 5.x хэлийг frontend болон backend аль аль дээр ашиглана. TypeScript нь
JavaScript-ийн дээр статик шалгалтыг нэмж, том хэмжээний системийн засвар
үйлчилгээг хялбарчилдаг.

\textbf{Өгөгдлийн сангийн удирдах системийн сонголт}

PostgreSQL 16 хувилбарыг сонгосон. ACID шаардлагыг бүрэн хангах, JSON/JSONB
дэмжлэг, оролцогчдын дийлэнх мэддэг SQL синтакс зэрэг давуу талтай.

\textbf{Баримт бичгийн боловсруулалтын системийн сонголт}

React компонент дэлгэцийн зохиомжийг Figma-д зохиомжилсны дараа TypeScript дээр
хэрэгжүүлсэн. Баримт бичгийг XeLaTeX дээр бичсэн.

\subsection{Техникийн шаардлага}

\textbf{Оролт гаралтын төхөөрөмжүүдийн шаардлага}

\begin{table}[htbp]
\caption{Техникийн доод шаардлага}
\label{tab:hw}
\centering
\small
\begin{tabularx}{\textwidth}{|L|L|}
\hline
\rowcolor{TableHdr}
\textcolor{white}{\textbf{Бүрэлдэхүүн}} &
\textcolor{white}{\textbf{Доод шаардлага}} \\ \hline
Зочны төхөөрөмж & iOS 14+ / Android 10+, камер, Wi-Fi/4G \\ \hline
\rowcolor{LightBlue}
Ажилтны төхөөрөмж & 10 дюйм дэлгэцтэй таблет эсвэл 1366×768 дэлгэцтэй зөөврийн компьютер \\ \hline
Сервер & 2 CPU core, 4 GB RAM, 50 GB SSD, 100 Mbps сүлжээ \\ \hline
\rowcolor{LightBlue}
Сүлжээ & Ресторан доторх Wi-Fi (-70 dBm буюу түүнээс дээш дохио) \\ \hline
\end{tabularx}
\end{table}

\section{Юзкейс диаграм}

Системийн гурван үндсэн хэрэглэгчийн юзкейс диаграмийг Зураг~\ref{fig:usecase} дэлгэрэнгүй харуулав.

\begin{figure}[htbp]
\centering
\begin{tikzpicture}[
  actor/.style={font=\small},
  uc/.style={ellipse, draw, minimum width=3.2cm, minimum height=0.8cm,
    align=center, font=\scriptsize},
  sys/.style={rectangle, draw, dashed, minimum width=5cm,
    minimum height=12cm, align=center},
  arr/.style={-{Stealth}},
  inc/.style={dashed, -{Stealth}}
]

% System boundary
\draw[dashed] (0,-6.5) rectangle (9,6);
\node[font=\small\bfseries] at (4.5,5.7) {Системийн хязгаар};

% Actors
\node[actor] (guest) at (-2.5, 2)  {\textbf{Зочин}};
\node[actor] (cashier) at (-2.5,-2) {\textbf{Кассир}};
\node[actor] (manager) at (11.5, 2) {\textbf{Менежер}};

% Guest use cases
\node[uc] (uc1) at (4,5)   {QR код уншуулах};
\node[uc] (uc2) at (4,3.5) {Цэс харах};
\node[uc] (uc3) at (4,2)   {Захиалга өгөх};
\node[uc] (uc4) at (4,0.5) {Статус харах};

% Cashier use cases
\node[uc] (uc5) at (4,-1)  {Захиалга харах};
\node[uc] (uc6) at (4,-2.5){Тооцоо баталгаажуулах};
\node[uc] (uc7) at (4,-4)  {Нэгтгэл харах};

% Manager use cases (right side also sees some)
\node[uc] (uc8) at (4,-5.5){Цэс удирдах};

% Associations
\draw[arr] (guest) -- (uc1);
\draw[arr] (guest) -- (uc2);
\draw[arr] (guest) -- (uc3);
\draw[arr] (guest) -- (uc4);
\draw[arr] (cashier) -- (uc5);
\draw[arr] (cashier) -- (uc6);
\draw[arr] (cashier) -- (uc7);
\draw[arr] (manager) -- (uc5);
\draw[arr] (manager) -- (uc7);
\draw[arr] (manager) -- (uc8);

% Include
\draw[inc] (uc3) -- node[right,font=\tiny]{<<include>>} (uc2);
\end{tikzpicture}
\caption{Системийн юзкейс диаграм}
\label{fig:usecase}
\end{figure}

\section{Юзкейс диаграмын тодорхойлолт}

\textbf{ЮК-01: QR код уншуулах}

\begin{table}[htbp]
\caption{ЮК-01 QR код уншуулах тодорхойлолт}
\label{tab:uc01}
\centering
\small
\begin{tabularx}{\textwidth}{|L|L|}
\hline
\rowcolor{LightBlue}
Юзкейсийн нэр & QR код уншуулах \\ \hline
Оролцогч & Зочин \\ \hline
\rowcolor{LightBlue}
Өмнөх нөхцөл & Зочин ширээнд суусан байна \\ \hline
Үндсэн урсгал & 1. Зочин камер нээнэ; 2. QR уншуулна; 3. Цэс харагдана \\ \hline
\rowcolor{LightBlue}
Дараах нөхцөл & Цэс харагдаж байна \\ \hline
Онцгой тохиолдол & QR хүчингүй бол алдааны мэдэгдэл харагдана \\ \hline
\end{tabularx}
\end{table}

\textbf{ЮК-02: Захиалга өгөх}

\begin{table}[htbp]
\caption{ЮК-02 Захиалга өгөх тодорхойлолт}
\label{tab:uc02}
\centering
\small
\begin{tabularx}{\textwidth}{|L|L|}
\hline
\rowcolor{LightBlue}
Юзкейсийн нэр & Захиалга өгөх \\ \hline
Оролцогч & Зочин \\ \hline
\rowcolor{LightBlue}
Өмнөх нөхцөл & Цэс харагдаж байна \\ \hline
Үндсэн урсгал & 1. Хоол сонгоно; 2. Сагсанд нэмнэ; 3. Захиалга батална; 4. Сервер хадгална; 5. Ажилтанд мэдэгдэнэ \\ \hline
\rowcolor{LightBlue}
Дараах нөхцөл & Захиалга pending статустай үүснэ \\ \hline
Онцгой тохиолдол & Хоол байхгүй болсон бол мэдэгдэнэ \\ \hline
\end{tabularx}
\end{table}

\textbf{ЮК-03: Тооцоо баталгаажуулах}

\begin{table}[htbp]
\caption{ЮК-03 Тооцоо баталгаажуулах тодорхойлолт}
\label{tab:uc03}
\centering
\small
\begin{tabularx}{\textwidth}{|L|L|}
\hline
\rowcolor{LightBlue}
Юзкейсийн нэр & Тооцоо баталгаажуулах \\ \hline
Оролцогч & Кассир \\ \hline
\rowcolor{LightBlue}
Өмнөх нөхцөл & Захиалга served статустай \\ \hline
Үндсэн урсгал & 1. Кассир захиалгыг сонгоно; 2. Нийт дүнг баталгаажуулна; 3. Paid болгоно \\ \hline
\rowcolor{LightBlue}
Дараах нөхцөл & Захиалга paid статусанд шилжинэ \\ \hline
Онцгой тохиолдол & Захиалга already paid бол алдаа гарна \\ \hline
\end{tabularx}
\end{table}

\textbf{ЮК-04: Цэс удирдах}

\begin{table}[htbp]
\caption{ЮК-04 Цэс удирдах тодорхойлолт}
\label{tab:uc04}
\centering
\small
\begin{tabularx}{\textwidth}{|L|L|}
\hline
\rowcolor{LightBlue}
Юзкейсийн нэр & Цэс удирдах \\ \hline
Оролцогч & Менежер \\ \hline
\rowcolor{LightBlue}
Өмнөх нөхцөл & Менежер нэвтэрсэн байна \\ \hline
Үндсэн урсгал & 1. Цэс хэсэгт нэвтэрнэ; 2. Хоол нэмэх/засах/устгах; 3. Зураг оруулах; 4. Хадгална \\ \hline
\rowcolor{LightBlue}
Дараах нөхцөл & Цэс шинэчлэгдэнэ, зочдод харагдана \\ \hline
Онцгой тохиолдол & Давхардсан нэр бол алдаа гарна \\ \hline
\end{tabularx}
\end{table}
